\documentclass{article}
\usepackage[margin=2cm]{geometry}
\usepackage{amsmath}
\usepackage{graphicx}
\usepackage{booktabs}
\usepackage[utf8]{inputenc}
\usepackage{ragged2e}
\setlength{\emergencystretch}{3em}
\begin{document}
\sloppy

\noindent\hfill\includegraphics[height=1.8em]{assets/figure_header_logo.png}\par\noindent\hrule\vspace{4pt}

This motivates us to define a superpotential W as


\begin{equation}
\begin{array}{c c}{{W=\sqrt{S_{0}^{2}+S_{3}^{2}}}}\end{array}\qquad (1)
\end{equation}


Unfortunately, this superpotential does not allow one to write the BPS equations for both
\$o and \$3 as gradient flow equations. The reason for this failure is that the integrability
condition required to convert the BPS equation into a gradient flow form is not satisfied;
see e.g. Appendix C.2.1 of . We thus follow the strategy of to construct an approximate
superpotential. Our model consists of two consistent truncations that admit flat domain walls.
and an exact superpotential. These are the \$3 = 0, \$o  0 truncation and the \$o = 0, \$3 = 0
truncation. The corresponding flow equations are (we set n = -1 henceforth).


\begin{equation}
\phi^{0^{\prime}}=-8\,\partial_{\bar{\phi}^{0}}W|_{\phi^{3}=0}~ \phi^{3^{\prime}}=8\,\partial_{\phi^{3}}W|_{\phi^{0}=0}~
\end{equation}


respectively. In either truncation, the BPS equations for the warp factor and dilaton o can.
be put in the following form,.


\begin{equation}
f^{\prime}=2\,W\,\qquad\qquad\sigma^{\prime}=2\,\partial_{\sigma}W\,\qquad (3)
\end{equation}


An important fact is that, though the gradient flow equations of do not hold exactly in the.
full model with \$o  0, \$3  0, they do hold up to and including O(z5). Looking at the form
of the UV asymptotics of the scalar fields, one may expand the superpotential of keeping.
only terms contributing up to this order. This gives.


\begin{equation}
W=\frac{1}{2}+\frac{3}{4}\sigma^{2}+\frac{1}{16}(\phi^{0})^{2}-\frac{3}{16}(\phi^{0})^{2}+\frac{1}{192}(\phi^{0})^{4}-\frac{3}{16}(\phi^{0})^{2}\sigma+\ldots~ (4)
\end{equation}


where the dots represent terms of order O(z6). This is the approximate superpotential we
will use in what follows..


\subsection*{0.0.1Bogomolnyi trick}


We now use the Bogomolnyi trick to get the finite counterterms needed to preserve supersymmetry in the case of a flat domain wall. The central idea of the Bogomolnyi trick is that
for a BPS solution, the renormalized on-shell action must vanish. In order to make use of
this fact, we will first want to recast the on-shell action in a simpler form. To do so, we
begin by inserting. We find that.


\begin{equation}
\begin{array}{c c}{{\displaystyle{\mathcal{L}}=-\frac{1}{4}R-20W^{2}+2{\mathcal{L}}_{\mathrm{kin}}~~~~~}}&{{~ 0000000000000000000000000000000000000000000000000000000000000000000000000000000000000000000000000000000000000000000000000000000000000000000000000000000000000000000000000000000000000000000000000000000000000000000000000000000}}\end{array}
\end{equation}


where we've defined


\begin{equation}
\mathcal{L}_{\mathrm{Aim}}=(\sigma^{\prime})^{2}+\frac{1}{4}\left[-(\phi^{3^{\prime}})^{2}+\cos^{2}{\phi}^{3}(\phi^{9})^{2}\right]\qquad (6)
\end{equation}

\end{document}
